\documentclass[11pt]{article}
\usepackage[margin=1.05in]{geometry}
\usepackage{amsmath,amssymb,amsthm,mathtools}
\usepackage{hyperref}

\theoremstyle{plain}
\newtheorem{thm}{Theorem}[section]
\newtheorem{lem}[thm]{Lemma}
\newtheorem{prop}[thm]{Proposition}
\newtheorem{cor}[thm]{Corollary}
\theoremstyle{definition}
\newtheorem{defn}[thm]{Definition}
\newtheorem{rem}[thm]{Remark}
\newtheorem{obs}[thm]{Observation}

\newcommand{\Z}{\mathbb{Z}}
\newcommand{\Saff}{\widetilde{S}_{N}}
\newcommand{\Fw}{\widetilde{F}_{w}}

\title{The exchange move is a re-timing, not a hop:\\
a TASEP reading of Theorem 4.1, and the death of a conjecture}
\author{Clio}
\date{22 September 2026 (PROVE c2, Day 200)}

\begin{document}
\maketitle

\begin{center}\small
\begin{tabular}{ll}
\textbf{Author:} Clio &
\textbf{Date:} 22 September 2026 \\
\textbf{Recipient:} Robin Langer, cc Rick (grandpa-rick) &
\textbf{Session:} PROVE c2, Q227 / Q221 \\
\end{tabular}\\[2pt]
\textbf{Title:} The exchange move is a re-timing, not a hop.
\end{center}

\begin{abstract}
DREAM c1 of 2026-09-22 conjectured (conjecture (C)) that the exchange-move set of
\cite[Thm.~4.1]{ClioExch} is ``the family of local moves that a TASEP exclusion rule
allows'', and offered a cheap falsifiable test: that the number of exchange moves at a
length-additive pair $(S,T)$ should equal the number of available hops of the corresponding
TASEP state. \textbf{The test fails, and (C) is false as stated.} The count of available hops
is not even an invariant of $v=u_Su_T$, whereas the count of exchange moves is
(\cite[Thm.~3.1]{ClioExch}, localisation); so no such identity can hold. Numerically the two
agree on $2111$ of $9168$ realisable configurations for $N\le 7$.

What replaces it is sharper than what was conjectured, and is proved here. Under Lam's
nilCoxeter action \cite[l.1700--1721]{Lam2006} a length-additive pair $(S,T)$ acting on a
cylindric shape is exactly \emph{two time-steps of a TASEP on a ring of $N$ sites}, and $S$
and $T$ are the sets of \emph{arcs} crossed at step~2 and step~1 respectively. On that locus
we prove (Theorem~\ref{thm:main}) that Theorem 4.1's letter $e$ is \emph{always} the bottom
$m$ of its run --- never an interior letter --- so that every exchange move is
\[
  S\mapsto S\setminus\{m\},\qquad T\mapsto T\cup\{m\},
\]
and (Theorem~\ref{thm:retime}) that this advances the intermediate configuration by exactly
one TASEP hop: the particle crosses its first arc one time-step earlier. The endpoints of the
trajectory, and the multiset of arcs crossed, are untouched. \textbf{The exchange move is a
re-timing of a fixed trajectory, not a hop of a state}, and that is precisely why no
single-state hop count can measure it.

The locus is exactly the 321-avoiding one, and the dichotomy is sharp: off it, $e>m$ occurs in
$2460$ of $6270$ pairs and the move has no trajectory reading at all. We also report a
negative that bears on the same dream: the Morse--Schilling excess of
\cite{ClioCrystal} is \emph{not} confined off the cylindric locus ($574$ of the $1001$ excess
instances for $N\le7$ are 321-avoiding), so the strictness proved on 09-22 c1 is not explained
by the cylindric/non-cylindric split.
\end{abstract}

\tableofcontents

\section{What was asked}\label{sec:asked}

Let $N\ge3$ and let $\Saff$ be the affine symmetric group, with $S\subsetneq\Z/N$ identified
with its $N$-periodic lift and $u_S$ the unique cyclically decreasing element of support $S$,
$\ell(u_S)=|S|$ \cite[Lem.~2.1]{ClioExch}. Call $(S,T)$ \emph{length-additive} if
$\ell(u_Su_T)=|S|+|T|$, and write $v=u_Su_T$.

\begin{thm}[{\cite[Thm.~4.1]{ClioExch}}]\label{thm:41}
Let $(S,T)$ be length-additive with $|S|>|T|$. Then there is a run $[m,M]$ of $S$ with
$m\notin T$ and an $e\in[m,M]$ with $e+1\notin T$; for the first such $e$ in the order
$m,m+1,\dots,M$, putting $S'=S\setminus\{e\}$, $T'=T\cup\{m\}$, one has $u_{S'}u_{T'}=v$ and
$\ell(v)=|S'|+|T'|$.
\end{thm}

Call a run $[m,M]$ of $S$ \emph{usable} if $m\notin T$ and $j+1\notin T$ for some $j\in[m,M]$.
Theorem~\ref{thm:41} produces one move per usable run; write $\mathcal{E}(S,T)$ for the set of
letters $e$ so produced, so $|\mathcal{E}(S,T)|=\#\{\text{usable runs of }S\}$.

The conjecture under test, from
\texttt{memory/connections/2026-09-22-a-crystal-operator-is-a-canonical-particle-hop.md}:

\begin{quote}
\textbf{(C)} \emph{The Theorem 4.1 exchange move set is the family of local moves that a TASEP
exclusion rule allows; Morse--Schilling's $x$-pairing is the tie-break that selects one of
them. Testable prediction: $|\mathcal{E}(S,T)|$ counts the available hops of the corresponding
TASEP state.}
\end{quote}

The brief for this session required a dictionary before any count
(\texttt{dictionary-before-identification}), and required that a mismatch be reported rather
than patched. Both requirements are met below: \S\ref{sec:dict} builds the dictionary,
\S\ref{sec:main} proves what is true on it, and \S\ref{sec:kill} reports the failure.

\begin{rem}[One half of the test was dropped before it was run]
As recorded in the WAKE addendum to
\texttt{questions/2026-09-22-dream-Q221-Q227-ledger.md}, the second half of Q227 --- ``is
$|\bigcup_x\tilde e_1^{(x)}(uv)|$ at most $|X(uv)|$?'' --- follows from the definition alone,
since $\tilde e_1^{(x)}$ is a function of $x$ and the union is a union of at most $|X(uv)|$
singletons. It is not instrumented here.
\end{rem}

\section{The dictionary: beads, arcs, and two time-steps}\label{sec:dict}

\subsection{Lam's action is an exclusion rule}

\begin{defn}[bead configuration]
A \emph{bead configuration} on the ring $\Z/N$ is a set $A\subseteq\Z$ with $A+N=A$ and
$\emptyset\ne A/N\Z\subsetneq\Z/N$. We read a bead as a particle and a non-bead as an empty
site, and we freely identify $A$ with $A/N\Z$.
\end{defn}

A cylindric shape of type $(N,m)$ in the bead model of \cite[\S2]{ClioCyl} is a strictly
increasing $x:\Z\to\Z$ with $x_{i+m}=x_i+N$; its bead set is $A=\{x_i\}$, which is a bead
configuration in the above sense. Adding one box per period to row $j$ is $x_j\mapsto x_j+1$,
legal exactly when $x_j+1\notin A$. Lam's nilCoxeter action \cite[l.1700--1721]{Lam2006}
(``$u_i\cdot\lambda=s_i\cdot\lambda$ if that adds boxes, else $0$''), transported to bead
coordinates by the derivation recorded in \cite[\S6.1]{Clio321} --- which fixes the convention
without any freedom except the origin of the labelling --- reads:
\begin{equation}\label{eq:ui}
  u_i\cdot A=\begin{cases}(A\setminus\{i\})\cup\{i+1\} & \text{if } i\in A,\ i+1\notin A,\\
  0 & \text{otherwise.}\end{cases}
\end{equation}

\begin{obs}
Equation~\eqref{eq:ui} is verbatim the totally asymmetric exclusion rule: \emph{the particle at
site $i$ hops one step clockwise if and only if site $i+1$ is empty.} This is not an analogy
imported from physics; it is Lam's box-adding rule written in the coordinates in which a
cylindric shape is a particle configuration on a ring. That standard cylindric tableaux are in
bijection with sequences of TASEP states is the main theorem of \cite{Elizalde2507}, whose
abstract reads, verbatim, ``\emph{We establish bijections between three classes of
combinatorial objects \dots\ lattice walks in simplicial regions \dots, standard cylindric
tableaux \dots, and sequences of states in the totally asymmetric simple exclusion process.}''
\end{obs}

\begin{rem}[Which paper supplies what]
\cite{Elizalde2507} (and its FPSAC~2026 extended abstract \cite{ElizaldeFPSAC}) supplies the
\emph{name} and the wider context: its $\Delta_{d,\ell}$ picture, with
$x_i=\alpha_{i-1}-\alpha_i$ \cite[\S4.2]{ElizaldeFPSAC}, is the gap-vector form of the same
ring, and its property (D3) --- \emph{every cylindric shape covers and is covered by
$c(\rho)=\#\{i:\rho_{i-1}>\rho_i\}$ elements} \cite[p.~5]{ElizaldeFPSAC} --- is the statement
that the numbers of available forward and backward hops agree. Nothing in
\S\S\ref{sec:dict}--\ref{sec:retime} \emph{consumes} \cite{Elizalde2507}: the proofs run on
\eqref{eq:ui} alone. The bridge from Morse--Schilling's affine-Stanley setting to Elizalde's
cylindric-tableau setting, which the DREAM note flagged as an open gap, is Lam's
\texttt{thm:321} together with its converse, quoted in \S\ref{sec:locus}; it is closed here for
the exchange move, and only for the exchange move.
\end{rem}

\subsection{What a cyclically decreasing element does to the beads}

Write $\mathrm{runs}(S)$ for the set of maximal cyclic runs of $S\subsetneq\Z/N$, each lifted
to an interval $[m,M]\subseteq\Z$ with $M-m+1\le|S|\le N-1$. The cyclically decreasing reduced
word for $u_S$ reads each run downwards, $s_Ms_{M-1}\cdots s_m$, runs concatenated in any
order; in the product $u_Su_T$ the \emph{rightmost} letter acts first
(\cite[Alternative Definition 2, l.853--861]{Lam2006}: ``\emph{the $a_1$ factor acts first,
i.e.\ is the rightmost}''). Hence in $v=u_Su_T$ the factor $u_T$ acts before $u_S$.

\begin{lem}[block travel]\label{lem:travel}
Let $A$ be a bead configuration and $S\subsetneq\Z/N$. Then $u_S\cdot A\ne0$ if and only if
for every run $[m,M]$ of $S$
\begin{equation}\label{eq:trav}
  m\in A\quad\text{and}\quad [m+1,M+1]\cap A=\emptyset ,
\end{equation}
and in that case
\[
  u_S\cdot A=\Bigl(A\setminus\{m:[m,M]\in\mathrm{runs}(S)\}\Bigr)\cup\{M+1:[m,M]\in\mathrm{runs}(S)\}.
\]
That is: \emph{for each run, the single bead at its bottom $m$ travels to $M+1$, crossing the
arcs $m,m+1,\dots,M$, and the whole interval it crosses must be empty.}
\end{lem}

\begin{proof}
Distinct runs $[m,M]$, $[m',M']$ of $S$ have disjoint blocks $[m,M+1]$, $[m',M'+1]$ modulo $N$:
if $M+1=m'$ then $m'\notin S$ and $m'\in S$, absurd; and the blocks cannot overlap otherwise
because the runs are separated by gaps. The operators for distinct runs therefore act on
disjoint sets of sites and commute, so it suffices to treat one run.

For one run the word is $s_Ms_{M-1}\cdots s_m$, applied rightmost-first, i.e.\ we apply
$u_m,u_{m+1},\dots,u_M$ in that order. We claim by induction on $j=m-1,m,\dots,M$ that after
applying $u_m,\dots,u_j$ the configuration is $A$ with the bead at $m$ replaced by a bead at
$j+1$, and that this is nonzero exactly when $m\in A$ and $[m+1,j+1]\cap A=\emptyset$. The base
$j=m-1$ is the empty statement. For the step, $u_{j+1}$ requires a bead at $j+1$ --- present,
it is the travelling bead --- and no bead at $j+2$ in the \emph{current} configuration. The
current configuration differs from $A$ only in that $m$ was removed and $j+1$ added, and
$j+2\notin\{m,j+1\}$: indeed $j+2$ ranges over $[m+2,M+1]$ and $M+1-m\le|S|\le N-1$, so
$j+2\not\equiv m$. Hence the condition is $j+2\notin A$. Composing the conditions over
$j=m-1,\dots,M-1$ gives exactly \eqref{eq:trav}.
\end{proof}

\begin{defn}[arcs]\label{def:arcs}
Label the arc from site $i$ to site $i+1$ by $i$. By Lemma~\ref{lem:travel}, if $u_S\cdot A\ne0$
then $S$ is precisely the set of arcs crossed by the beads during the step $A\mapsto u_S\cdot A$,
each arc crossed exactly once, each particle crossing a contiguous block of arcs.
\end{defn}

So $\ell(u_S)=|S|$ is the total displacement of the step --- equivalently, by
\cite[\S2]{ClioCyl}, the number of cells of the horizontal strip added --- and the
\emph{trajectory} associated with a length-additive pair is:
\begin{equation}\label{eq:traj}
  A\ \xrightarrow{\ \text{step 1, arcs }T\ }\ \nu:=u_T\cdot A
     \ \xrightarrow{\ \text{step 2, arcs }S\ }\ \lambda:=u_S\cdot\nu .
\end{equation}

\subsection{The locus on which the dictionary exists}\label{sec:locus}

\begin{defn}
A length-additive pair $(S,T)$ is \emph{realisable} if there is a bead configuration $A$ with
$u_S\cdot(u_T\cdot A)\ne0$, i.e.\ if the trajectory \eqref{eq:traj} exists.
\end{defn}

\begin{prop}[the locus is the 321-avoiding one]\label{prop:locus}
$(S,T)$ is realisable if and only if $v=u_Su_T$ is 321-avoiding.
\end{prop}

\begin{proof}
($\Leftarrow$) is \cite[\texttt{thm:321}, l.1844--47]{Lam2006} in the form recorded in
\cite[Thm.~5.2]{Clio321}. ($\Rightarrow$) is Lam's converse,
\cite[l.1871--73]{Lam2006}, verbatim: \emph{``If $\lambda$ and $\mu$ are cylindric shapes
satisfying $u_w\cdot\mu=\lambda$ then $w$ is necessarily 321-avoiding.''}
Independently verified here on all $10729$ length-additive pairs for $3\le N\le7$: agreement
$10729/10729$ (\S\ref{sec:verify}, \textbf{V2}).
\end{proof}

\section{The main theorem: the letter is always the run bottom}\label{sec:main}

\begin{thm}\label{thm:main}
Let $A$ be a bead configuration, let $(S,T)$ be length-additive, and suppose
$\nu=u_T\cdot A\ne0$ and $\lambda=u_S\cdot\nu\ne0$. Let $[m,M]$ be any run of $S$. Then
\begin{enumerate}
\item[(a)] $m\notin T$;
\item[(b)] either $m+1\notin T$, or $[m+1,M+1]\subseteq T$;
\item[(c)] the run is usable if and only if $m+1\notin T$, and in that case Theorem~\ref{thm:41}'s
letter is $e=m$.
\end{enumerate}
Consequently every exchange move of Theorem~\ref{thm:41} at such an $(S,T)$ is
\[
  S\ \mapsto\ S\setminus\{m\},\qquad T\ \mapsto\ T\cup\{m\},
\]
for a run bottom $m$ of $S$ with $m+1\notin T$, and
\begin{equation}\label{eq:count}
  |\mathcal{E}(S,T)|\;=\;\#\{[m,M]\in\mathrm{runs}(S)\ :\ m+1\notin T\}.
\end{equation}
\end{thm}

\begin{proof}
Throughout, Lemma~\ref{lem:travel} is applied twice: to $T$ acting on $A$, and to $S$ acting on
$\nu$. From the second,
\begin{equation}\label{eq:star}
  m\in\nu\quad\text{and}\quad[m+1,M+1]\cap\nu=\emptyset. \tag{$*$}
\end{equation}
From the first, $\nu=\bigl(A\setminus\{\text{run bottoms of }T\}\bigr)\cup\{b+1:[a,b]\in
\mathrm{runs}(T)\}$, and each run $[a,b]$ of $T$ satisfies $a\in A$, $[a+1,b+1]\cap A=\emptyset$.

\emph{(a).} Suppose $m\in T$ and let $[a,b]$ be the run of $T$ containing $m$.
If $m=a$, then the bead at $a$ left, and $a$ receives a bead in $\nu$ only if $a=b'+1$ for some
run $[a',b']$ of $T$; but $b'+1\notin T$ while $a\in T$. So $m=a\notin\nu$, contradicting
\eqref{eq:star}. If $a<m\le b$, then $m\in[a+1,b+1]$ so $m\notin A$, and again $m\in\nu$ would
force $m=b'+1\notin T$, contradicting $m\in T$. So $m\notin\nu$, again contradicting
\eqref{eq:star}. Hence $m\notin T$.

\emph{(b).} Suppose $m+1\in T$. By (a), $m\notin T$, so the run of $T$ containing $m+1$ begins
at $m+1$: call it $[m+1,b]$, with $b\le m+N-1$. By Lemma~\ref{lem:travel} the bead travelling
along it lands at $b+1$, so $b+1\in\nu$. By \eqref{eq:star}, $b+1\notin[m+1,M+1]$; since
$b+1\ge m+2$, this forces $b+1\ge M+2$, i.e.\ $b\ge M+1$. Hence
$[m+1,M+1]\subseteq[m+1,b]\subseteq T$.

\emph{(c).} By (a) the run passes the first test of usability. If $m+1\notin T$, then $j=m$
satisfies $j+1\notin T$, so the run is usable and the \emph{first} such $j$ is $e=m$. If
$m+1\in T$, then by (b) $[m+1,M+1]\subseteq T$, so $j+1\in T$ for every $j\in[m,M]$ and the run
is not usable. Equation~\eqref{eq:count} follows, since Theorem~\ref{thm:41} contributes exactly
one move per usable run.
\end{proof}

\begin{rem}[why an interior letter is impossible, in one sentence]
In an exclusion process only the \emph{rear} particle of a moving block can be held back: if an
interior arc of a particle's journey were re-timed, the particle behind it would have to pass
through it. Theorem~\ref{thm:main}(c) is that sentence, proved.
\end{rem}

\begin{rem}[the hypothesis $|S|>|T|$ is not used]
Theorem~\ref{thm:main} places no hypothesis relating $|S|$ and $|T|$. In
\cite[Thm.~4.1]{ClioExch} the hypothesis $|S|>|T|$ is consumed \emph{only} by the
applicability lemma, which guarantees that a usable run \emph{exists}; the move itself needs
only a usable run. Comparing against what the theorem guarantees rather than what it
constructs is the error recorded in
\texttt{a-hypothesis-of-my-own-theorem-is-an-unexamined-assumption}, and it is avoided here.
\end{rem}

\section{The move is a re-timing}\label{sec:retime}

\begin{thm}\label{thm:retime}
In the situation of Theorem~\ref{thm:main}, let $[m,M]$ be a run of $S$ with $m+1\notin T$, and
put $T'=T\cup\{m\}$, $S'=S\setminus\{m\}$. Then
\[
  \nu':=u_{T'}\cdot A=(\nu\setminus\{m\})\cup\{m+1\}\ne0,
  \qquad u_{S'}\cdot\nu'=\lambda .
\]
That is: the exchange move advances the intermediate configuration by exactly one TASEP hop ---
of the very particle that is about to make the step-2 journey $m\to M+1$ --- and leaves the
endpoints $A$ and $\lambda$ of the trajectory unchanged. Equivalently, in the language of
Definition~\ref{def:arcs}, it moves the crossing of the single arc $m$ from step 2 to step 1.
\end{thm}

\begin{proof}
By Theorem~\ref{thm:main}(a),(c) we have $m\notin T$ and $m+1\notin T$. Two facts first.

\emph{$m\in A$ or $m=b+1$ for the run of $T$ ending at $m-1$.} Since $m\in\nu$ by
\eqref{eq:star}, either $m\in A$ and $m$ is not a run bottom of $T$ (which holds, as $m\notin T$),
or $m=b+1$ for some run $[a,b]$ of $T$, i.e.\ $b=m-1\in T$.

\emph{$m+1\notin A$.} By \eqref{eq:star}, $m+1\notin\nu$. If $m+1\in A$ then, $\nu$ being
obtained from $A$ by deleting run bottoms of $T$ and adding tops-plus-one, $m+1$ can be absent
from $\nu$ only if $m+1$ is a run bottom of $T$; but $m+1\notin T$. So $m+1\notin A$.

\emph{Case $m-1\notin T$.} Then $\{m\}$ is a singleton run of $T'$ and every other run of $T'$
is a run of $T$. By the first fact $m\in A$; by the second $m+1\notin A$; so
condition~\eqref{eq:trav} holds for the run $\{m\}$. The blocks of the other runs are unchanged
and are disjoint from $[m,m+1]$: a block $[a,b+1]$ contains $m$ only if $m\in T$ (no) or
$m=b+1$, i.e.\ $m-1=b\in T$ (excluded in this case); and it contains $m+1$ only if $m+1\in T$
(no) or $m+1=b+1$, i.e.\ $m\in T$ (no). By Lemma~\ref{lem:travel},
$\nu'=(\nu\setminus\{m\})\cup\{m+1\}$.

\emph{Case $m-1\in T$.} Let $[a,m-1]$ be the run of $T$ ending at $m-1$; the corresponding run
of $T'$ is $[a,m]$. Condition~\eqref{eq:trav} for $T$ gives $a\in A$ and
$[a+1,m]\cap A=\emptyset$; with $m+1\notin A$ this is exactly \eqref{eq:trav} for $[a,m]$. All
other runs are untouched. By Lemma~\ref{lem:travel} the travelling bead now lands at $m+1$
instead of $m$, so again $\nu'=(\nu\setminus\{m\})\cup\{m+1\}$.

\emph{The second step.} $S'=S\setminus\{m\}$ replaces the run $[m,M]$ by $[m+1,M]$ (empty if
$M=m$) and leaves the other runs of $S$ alone. Condition~\eqref{eq:trav} for $[m+1,M]$ on $\nu'$
reads $m+1\in\nu'$ --- true by construction --- and $[m+2,M+1]\cap\nu'=\emptyset$, which follows
from \eqref{eq:star} because $\nu'$ differs from $\nu$ only at $m$ (removed) and $m+1$ (added),
neither of which lies in $[m+2,M+1]$. The other runs of $S$ have blocks disjoint from
$[m,M+1]$, so their conditions on $\nu'$ are the same as on $\nu$. By Lemma~\ref{lem:travel} the
bead now at $m+1$ travels to $M+1$, and every other bead does what it did before; hence
$u_{S'}\cdot\nu'=u_S\cdot\nu=\lambda$.
\end{proof}

\begin{cor}[a second proof of Theorem~\ref{thm:41} on the cylindric locus]\label{cor:second}
For realisable $(S,T)$, Theorem~\ref{thm:retime} re-proves the identity $u_{S'}u_{T'}=v$ of
\cite[Lem.~4.3]{ClioExch} by particle bookkeeping alone, with no conjugation computation in
$\Saff$: the two factorisations move the same beads along the same arcs and differ only in when
the arc $m$ is crossed. The original proof is the one that covers the whole of $\Saff$; this
one covers only the 321-avoiding locus, and is included because it is the reason the statement
is true there.
\end{cor}

\begin{cor}[the arc multiset is what $v$ remembers]\label{cor:multiset}
The multiset $S\uplus T$ of arcs crossed is invariant under the exchange move. Computationally
(\S\ref{sec:verify}, \textbf{V5}) it is constant on the whole fibre $\mathcal{F}(v)$ exactly
when $v$ is 321-avoiding: constant on all $1343$ 321-avoiding fibres for $N\le7$, and on no
other fibre. \emph{This is reported as computed, not proved.}
\end{cor}

\section{The death of (C)}\label{sec:kill}

\begin{thm}[(C) is false]\label{thm:kill}
There is no identity between $|\mathcal{E}(S,T)|$ and the number of available hops of the state
$\nu$ of \eqref{eq:traj}, for two independent reasons.
\begin{enumerate}
\item[\textbf{(K1)}] \emph{Structural.} By \cite[Thm.~3.1]{ClioExch} (localisation),
$\mathcal{E}(S,T)$ --- indeed the whole fibre --- depends only on $v=u_Su_T$. The number of
available hops of $\nu$ does not: it depends on the base point $A$. Smallest non-degenerate
witness, $N=4$, $S=\{0\}$, $T=\{1\}$:
\[
  A=\{0,1\}\Rightarrow\nu=\{0,2\}\ \text{with 2 available hops},\qquad
  A=\{0,1,3\}\Rightarrow\nu=\{0,2,3\}\ \text{with 1}.
\]
A quantity that is not a function of $v$ cannot equal one that is. Non-constancy occurs for
$1166$ of the $4459$ realisable pairs with $N\le7$.
\item[\textbf{(K2)}] \emph{Numerical.} Over the $9168$ realisable triples $(S,T,A)$ with
$N\le7$, $|\mathcal{E}(S,T)|$ equals the number of available forward hops of $\nu$ in $2111$
cases; of $\lambda$ in $2258$; of $A$ in $2258$. The smallest non-degenerate failure is $N=3$,
$S=\{0\}$, $T=\{1\}$, $A=\{0,1\}$: no exchange move, one available hop.
\end{enumerate}
By contrast the replacement count \eqref{eq:count} holds in $9168$ of $9168$ cases.
\end{thm}

\begin{proof}
(K1) is Theorem~\ref{thm:41}'s localisation together with the displayed witness, which is
checked directly from \eqref{eq:ui}. (K2) is a finite computation, \S\ref{sec:verify}.
\end{proof}

\subsection{What survives, and what the right statement is}

The mechanism the dream identified is real; the estimand was wrong. The exchange move set is
not a set of hops of a \emph{state}. It is a set of re-timings of a two-step \emph{trajectory}
whose endpoints, particles and total displacements are already fixed by $v$. This is exactly
why a single-state hop count cannot measure it, and it is visible in
Corollary~\ref{cor:multiset}: what $v$ determines is the multiset of arcs to be crossed; what
the fibre $\mathcal{F}(v)$ parametrises is the schedule.

\begin{quote}
\textbf{(C$'$), proved (Theorems~\ref{thm:main},~\ref{thm:retime}).} On the 321-avoiding locus,
the Theorem 4.1 exchange moves at $(S,T)$ are in bijection with the step-2 journeys whose first
two arcs are both free at step 1, and each move advances the intermediate configuration by one
TASEP hop of that journey's particle. The movable arc is always the \emph{first} arc of a
journey, never an interior one, because in an exclusion process only the rear of a moving block
can be re-timed.
\end{quote}

\begin{rem}[the locality claim, made precise]
(C) called the Theorem 4.1 condition ``local''. Theorem~\ref{thm:main} says how local: the
usability of a run $[m,M]$ is decided by the two sites $m$, $m+1$ alone
(``$m\notin T$ and $m+1\notin T$''), whereas off the locus the condition genuinely scans the
run, and $e>m$ occurs in $2460$ of the $6270$ non-321-avoiding pairs with $N\le7$. So locality
is not a property of Theorem 4.1; it is a property of Theorem 4.1 \emph{restricted to the
cylindric locus}, and the restriction is what buys it.
\end{rem}

\subsection{A second negative: the Morse--Schilling excess is not the locus}

\cite{ClioCrystal} proved that $\bigcup_{x}\{\tilde e_1^{(x)}(uv)\}\subsetneq\mathcal{E}(S,T)$,
strictly, first at $N=5$, with excess $0,0,15,132,854$ for $N=3,\dots,7$ (reproduced here as a
positive control, \S\ref{sec:verify}, \textbf{V1}). It would have been attractive if the excess
were exactly the non-cylindric part. It is not: of the $1001$ instances with excess,
$574$ are 321-avoiding (by $N$: $0,0,15,90,469$). So the strictness of the Morse--Schilling
containment is \emph{not} explained by the cylindric/non-cylindric split, and the dream's
suggestion that the strictness ``is the price of being a function'' is neither confirmed nor
refuted here. It remains open.

\section{Computational verification}\label{sec:verify}

Code at \texttt{proofs/code-q227/} (\texttt{beads.py}, \texttt{order.py}, \texttt{lemma1.py},
\texttt{verdict.py}, \texttt{kill.py}, \texttt{invar.py}), running on the existing
\texttt{proofs/code-q220/} harness. Ranges: all $10729$ length-additive pairs for
$3\le N\le7$ (the $9093$-instance census of \cite{ClioCrystal} is these minus the $1636$ pairs
with $X(uv)=\emptyset$), and all $9168$ realisable triples $(S,T,A)$.

\begin{itemize}
\item[\textbf{V0}] \emph{Positive control, inherited.} \texttt{code-q220/final.py} reproduces
the published Morse--Schilling excess sequence $0,0,15,132,854$ and the pair counts
$34,143,561,2122,7869$. The Morse--Schilling implementation it rests on reproduces all three
worked examples printed in \cite{MorseSchilling}; nothing is fitted.
\item[\textbf{V1}] \emph{Lemma~\ref{lem:travel}, exhaustively.} Predicted nonzero-ness
\eqref{eq:trav} and predicted image checked against the composed action \eqref{eq:ui} on all
$85860$ pairs $(S,A)$ with $3\le N\le8$: $85860$ agreements, $0$ disagreements. \emph{Could
fail:} an off-by-one in the operator order or in the wrap-around would break it immediately.
\item[\textbf{V2}] \emph{Proposition~\ref{prop:locus}.} ``realisable'' versus ``321-avoiding''
(tested by the pattern criterion on the window over three periods): agreement
$10729/10729$.
\item[\textbf{V3}] \emph{Theorem~\ref{thm:main}.} On all $4459$ realisable pairs: $m\notin T$
for every run of $S$, $4459/4459$; $e=m$ for every usable run, $4459/4459$; every unusable run
satisfies $[m+1,M+1]\subseteq T$, $4459/4459$; and the count \eqref{eq:count},
$4459/4459$ pairs and $9168/9168$ triples. Off the locus, $e>m$ in $2460$ of $6270$ pairs ---
so the claim has real content and is not an artefact of the definition.
\item[\textbf{V4}] \emph{Theorem~\ref{thm:retime} and Corollary~\ref{cor:second}.} For every
realisable triple and every exchange move, the image of the full word for $u_{S'}u_{T'}$ on $A$
equals that of $u_Su_T$: $1326/1326$ at $N=6$ and likewise below, $0$ exceptions.
\item[\textbf{V5}] \emph{Corollary~\ref{cor:multiset}.} $S\uplus T$ constant on
$\mathcal{F}(v)$: $1343/1343$ 321-avoiding fibres, and $1343$ of $3019$ fibres in total, the
two counts agreeing at every $N$ --- so constancy holds exactly on the 321-avoiding fibres.
\item[\textbf{V6}] \emph{Theorem~\ref{thm:kill}.} As displayed in \S\ref{sec:kill}.
\end{itemize}

\begin{rem}[what the computation could not have caught]
Every observable above is computed from \eqref{eq:ui}. If the transport of Lam's rule into bead
coordinates were wrong, \textbf{V2} would be the detector: realisability would not track
321-avoidance, and Lam's dichotomy is sharp enough that a wrong convention would break it. It
does not break. This is the same control used in \cite[\S6]{Clio321}, run again here
independently.
\end{rem}

\section{Gaps}\label{sec:gaps}

\begin{enumerate}
\item \textbf{Corollary~\ref{cor:multiset} is computed, not proved.} The forward direction
(constant on 321-avoiding fibres) should follow from Lemma~\ref{lem:travel} plus the uniqueness
statement \cite[l.1723]{Lam2006} (``at most one $w$ with $u_w\cdot\mu=\lambda$''), since
$S\uplus T$ is then the content multiset of $\lambda/A$; the converse is not attempted.
\item \textbf{The Morse--Schilling excess is unexplained.} \S\ref{sec:kill} rules out the
cylindric split as the explanation. The dream's claim that the strictness is the price of
canonicity is untested.
\item \textbf{Semistandard versus standard, and the growth-diagram route.} Nothing here uses
\cite{ElizaldeFPSAC}'s local rules (F1)/(F2) or the correspondence CRSK. The connection
asserted in \S\ref{sec:dict} is at the level of the underlying particle model only. A genuine
statement relating the exchange move to cylindric growth diagrams would need the rules
themselves, and is not claimed.
\item \textbf{Novelty.} No claim of novelty is made for the TASEP reading of cylindric
tableaux, which is \cite{Elizalde2507}, nor for the 321-avoiding $\leftrightarrow$ cylindric
correspondence, which is \cite{Lam2006}. What is claimed new is the \emph{route}: the
identification of the exchange move of \cite{ClioExch} with a one-arc re-timing, and the
consequent proof that $e=m$ on the locus.
\end{enumerate}

\begin{thebibliography}{ClioCrystal}
\bibitem[ClioCyl]{ClioCyl} Clio, \emph{M-convexity of cylindric skew Schur polynomials},
\texttt{proofs/2026-09-20-c1-cylindric-M-convexity.tex}, 20 September 2026.
\bibitem[ClioExch]{ClioExch} Clio, \emph{Localisation and an exchange move for affine Stanley
symmetric polynomials}, \texttt{proofs/2026-09-21-c1-affine-stanley-exchange.tex},
21 September 2026.
\bibitem[Clio321]{Clio321} Clio, \emph{M-convexity and the Newton polytope of affine Stanley
symmetric polynomials on the 321-avoiding class},
\texttt{proofs/2026-09-21-c2-affine-stanley-321-avoiding.tex}, 21 September 2026.
\bibitem[ClioCrystal]{ClioCrystal} Clio, \emph{Crystal edges are exchange moves},
\texttt{proofs/2026-09-22-c1-crystal-edges-are-exchange-moves.tex}, 22 September 2026.
\bibitem[Eliz-A]{ElizaldeFPSAC} S. Elizalde, \emph{Cylindric growth diagrams},
S\'em. Lothar. Combin. XX (2026), Article \#26, 12 pp. (FPSAC 2026, Seattle).
\bibitem[Eliz-B]{Elizalde2507} S. Elizalde, \emph{Cylindric growth diagrams, walks in
simplices, and exclusion processes}, arXiv:2507.01097; Electron. J. Combin. 33(3) (2026) \#P7.
\bibitem[Lam]{Lam2006} T. Lam, \emph{Affine Stanley symmetric functions},
Amer. J. Math. 128 (2006) 1553--1586; arXiv:math/0501335. \emph{Cited by internal label and
line number of the arXiv source, never by printed number, per
\texttt{a-derived-identifier-is-not-a-quotation}.}
\bibitem[MS]{MorseSchilling} J. Morse and A. Schilling, \emph{A crystal approach to affine
Schubert calculus}, arXiv:1408.0320. \emph{Cited by internal label and line number.}
\end{thebibliography}

\end{document}
